\section*{第 \ref{ch:g8:negprod} 章——负数的乘法}

\begin{solution}{exo:g8:negprod:1}
$(-8) \times (+7) = -56$; \quad
$(-9) \times (-6) = +54$; \quad
$(+12) \times (-0.5) = -6$; \quad
$(-1) \times (-1) \times (-1) = -1$（三个负\omterm{def:g5:division:multiple}{因数}，\omterm{def:g3:numbers:evenodd}{奇数}）。
\end{solution}

\begin{solution}{exo:g8:negprod:2}
$(-63) \div (+9) = -7$; \quad
$(-8) \div (-16) = +0.5$; \quad
$\dfrac{45}{-5} = -9$; \quad
$\dfrac{-7.2}{-8} = +0.9$。
\end{solution}

\begin{solution}{exo:g8:negprod:3}
三个负\omterm{def:g5:division:multiple}{因数}：\emph{负}。\quad
$(-1)^{10}$：十个负\omterm{def:g5:division:multiple}{因数}，\omterm{def:g3:numbers:evenodd}{偶数}个，\emph{正}（其值为
$1$）。\quad
$(-2)^7$：七个\omterm{def:g5:division:multiple}{因数}，\omterm{def:g3:numbers:evenodd}{奇数}个，\emph{负}。\quad
$(-4) \times 0 \times (-6) = 0$：既非正数也非负数。
\end{solution}

\begin{solution}{exo:g8:negprod:4}
$(-2)^4 = +16$；\quad $-2^4 = -16$（负号不参与乘方）；\quad
$(-3)^3 = -27$；\quad $(-1)^{2026} = +1$（指数为\omterm{def:g3:numbers:evenodd}{偶数}）。
\end{solution}

\begin{solution}{exo:g8:negprod:5}
$7 + 2 \times (-6) = 7 + (-12) = -5$。

$(-4) \times (5 - 9) = (-4) \times (-4) = +16$。

$18 \div (-3) - (-2) = -6 + 2 = -4$。
\end{solution}

\begin{solution}{exo:g8:negprod:6}
$\dfrac{(-5) \times (+8)}{-4} = \dfrac{-40}{-4} = +10$。

$\dfrac{(-3) \times (-6)}{(-2) \times (+9)} = \dfrac{+18}{-18} = -1$。
\end{solution}

\begin{solution}{exo:g8:negprod:7}
当 $x = -3$ 时：\ $5x = -15$；\quad $x^2 = (-3)^2 = 9$；\quad
$-x^2 = -9$; \quad $2x^2 + 4x = 18 - 12 = 6$; \quad
$(2x)^2 = (-6)^2 = 36$。
\end{solution}

\begin{solution}{exo:g8:negprod:8}
$(-6) \times (-7) = 42$; \quad
$10 \div (-4) = -2.5$; \quad
$(-5) \times \frac15 = -1$，所以缺的数是
$+\frac15 = 0.2$。
\end{solution}

\begin{solution}{exo:g8:negprod:9}
\emph{1. 错误}：$0.5 \times 0.5 = 0.25$ 小于
$0.5 + 0.5 = 1$（或：$2 \times 3 = 6 > 5$，但结论须
对一切情形都成立）。

\emph{2. 正确}：一个数的平方是两个同号数的积，
因此为正数或零。

\emph{3. 错误}：$(-2) \times (-3) \times 4 = +24$ 为正数，
其中有两个负\omterm{def:g5:division:multiple}{因数}。
\end{solution}

\begin{solution}{exo:g8:negprod:10}
设 $r$ 为答对的题数；则答错的有 $20 - r$ 题，且
\[
5r - 3(20 - r) = 36
\ \Longrightarrow\
5r - 60 + 3r = 36
\ \Longrightarrow\
8r = 96
\ \Longrightarrow\
r = 12 .
\]
答对 12 题，答错 8 题。验算：
$60 - 24 = 36$。
\end{solution}

\begin{solution}{exo:g8:negprod:11}
零乘任何数都得零，所以 $(-1) \times 0 = 0$。
把 $0 = 1 + (-1)$ 代入并用分配律展开：
\[
0 = (-1) \times \bigl(1 + (-1)\bigr)
= (-1) \times 1 + (-1) \times (-1)
= -1 + (-1) \times (-1).
\]
数 $(-1) \times (-1)$ 与 $-1$ 相加得 $0$：它必须是
$-1$ 的\omterm{def:g7:negatives:def}{相反数}，即 $+1$。由分配律可知，
别无其他选择。
\end{solution}

\begin{solution}{pb:g8:negprod:1}
\textbf{1.} 由 $0 = 0 + 0$，利用分配律
（\cref{thm:g7:priorities:distributivity}）得
\[
0 \times a = (0 + 0) \times a = 0 \times a + 0 \times a .
\]
于是，$0 \times a$ 与自身相加结果不变，而加上后不改变任何数的数只有
$0$。（直接地：
两边同减 $0 \times a$。）因此 $0 \times a = 0$。

\textbf{2.} 展开后，利用 $1 + (-1) = 0$ 及第 1 问的结论：
\[
a + (-1) \times a
= 1 \times a + (-1) \times a
= \bigl(1 + (-1)\bigr) \times a
= 0 \times a = 0 .
\]
因此 $(-1) \times a$ 是与 $a$ 相加得 $0$ 的数：
即 $a$ 的\omterm{def:g7:negatives:def}{相反数}，故 $(-1) \times a = -a$。
\Cref{exo:g8:negprod:11} 是 $a = -1$ 时的特殊情形，此时
$(-1) \times (-1) = -(-1) = +1$。

\textbf{3.} 记 $-a = (-1) \times a$（第 2 问），并重新组合
各因子：
\[
(-a) \times b
= \bigl((-1) \times a\bigr) \times b
= (-1) \times (a \times b)
= -(a \times b) .
\]

\textbf{4.} 对两个因子各应用一次第 3 问：
\[
(-a) \times (-b)
= -\bigl(a \times (-b)\bigr)
= -\bigl(-(a \times b)\bigr)
= a \times b ,
\]
因为一个数的\omterm{def:g7:negatives:def}{相反数}的\omterm{def:g7:negatives:def}{相反数}就是这个数本身。

\textbf{5.} 记 $a = +d$ 或 $a = -d$，$b = +e$ 或 $b = -e$，
其中 $d$ 与 $e$ 分别为两数到零的距离。四种情形为：
\[
(+d) \times (+e) = de, \quad
(+d) \times (-e) = -de, \quad
(-d) \times (+e) = -de, \quad
(-d) \times (-e) = +de,
\]
中间两式由第 3 问得出，末式由第 4 问得出。各情形下，
结果的距离为 $d \times e$，即两距离的积，
符号由
\cref{thm:g8:negprod:rules}的表格给出：同号得正，
异号得负。

\textbf{6.} $(-1)^2 = +1$，$(-1)^3 = -1$，$(-1)^4 = +1$，
$(-1)^5 = -1$。一般地，$(-1)^n$ 是 $n$ 个负因子的积，
每个因子的距离均为 $1$。由\cref{prop:g8:negprod:many}可知，
其值为 $+1$（$n$ 为\omterm{def:g3:numbers:evenodd}{偶数}时）或 $-1$（$n$ 为\omterm{def:g3:numbers:evenodd}{奇数}时）。

\textbf{7.} 共有十个负因子——\omterm{def:g3:numbers:evenodd}{偶数}个——故
积为\emph{正数}（\cref{prop:g8:negprod:many}）。其
距离为 $1 \times 2 \times 3 \times \dots \times 10$。

\textbf{8.} 由第 4 问，$(-a)^2 = (-a) \times (-a) = a \times a =
a^2$。再由第 3 问，
\[
(-a)^3 = (-a)^2 \times (-a) = a^2 \times (-a) = -(a^2 \times a)
= -a^3 .
\]
负号在指数为\omterm{def:g3:numbers:evenodd}{偶数}时消失，在指数为\omterm{def:g3:numbers:evenodd}{奇数}时保留，
与\cref{ex:g8:negprod:powers}一致。

\textbf{9.} 若 $a = 0$，则 $a^2 = 0$。否则 $a$ 与 $a$ 具有
\emph{同号}，故 $a^2 = a \times a$ 为正数
（\cref{thm:g8:negprod:rules}）。因此一个数的平方永不为
负；由于 $-4$ 为负数，故没有任何\omterm{def:g7:negatives:def}{正负数} $x$ 满足
$x^2 = -4$。

\textbf{10.} 佐伊说得对。指数 $2026$ 与 $2027$ 的奇偶性
相反，故 $(-1)^{2026} = +1$，$(-1)^{2027} = -1$，
两者之和为 $0$。一般地，相邻两个指数中一个为\omterm{def:g3:numbers:evenodd}{偶数}，
一个为\omterm{def:g3:numbers:evenodd}{奇数}，故这两个幂为 $+1$ 与 $-1$（次序不定），
即互为\omterm{def:g7:negatives:def}{相反数}，其和总为 $0$。

\textbf{11.} $A_3 = 1 - 2 + 3 = 2$，再
$A_4 = A_3 - 4 = -2$, $A_5 = A_4 + 5 = 3$, $A_6 = A_5 - 6 = -3$,
$A_7 = A_6 + 7 = 4$。猜想：$A_n = -\frac{n}{2}$（$n$ 为\omterm{def:g3:numbers:evenodd}{偶数}），
$A_n = \frac{n+1}{2}$（$n$ 为\omterm{def:g3:numbers:evenodd}{奇数}）。

\textbf{12.} 当 $n$ 为\omterm{def:g3:numbers:evenodd}{偶数}时，各项恰好两两配对：
\[
A_n = (1 - 2) + (3 - 4) + \dots + \bigl((n-1) - n\bigr),
\]
每对之和均为 $-1$。共有 $\frac{n}{2}$ 对（每对两项，
共 $n$ 项），于是
$A_n = \frac{n}{2} \times (-1) = -\frac{n}{2}$.

\textbf{13.} 当 $n$ 为\omterm{def:g3:numbers:evenodd}{奇数}时，和 $A_n$ 为 $A_{n-1}$（前 $n - 1$ 项之和）
加上最后一项 $+n$，因 $n$ 为\omterm{def:g3:numbers:evenodd}{奇数}。
由于 $n - 1$ 为\omterm{def:g3:numbers:evenodd}{偶数}，由第 12 问得
\[
A_n = A_{n-1} + n = -\frac{n-1}{2} + n
= \frac{-(n-1) + 2n}{2} = \frac{n+1}{2} .
\]

\textbf{14.} $A_{100} = -\frac{100}{2} = -50$, \quad
$A_{101} = \frac{102}{2} = 51$, \quad
$A_{2026} = -\frac{2026}{2} = -1013$.

\textbf{15.} 并不意外：由 $A_{100}$ 到 $A_{101}$ 恰好增加
一项 $+101$，故其差必为 $101$。
同理，$A_{2027} - A_{2026} = +2027$——无需
计算 $A_{2027}$ 本身。
\end{solution}
